\documentclass[11pt]{article}
\usepackage[margin=1.1in]{geometry}
\usepackage{amsmath,amssymb,amsthm,mathtools}
\usepackage{enumitem}
\usepackage[hidelinks]{hyperref}
\expandafter\def\expandafter\UrlBreaks\expandafter{\UrlBreaks\do\-}
\usepackage{microtype}
\emergencystretch=3em

\newtheorem{theorem}{Theorem}[section]
\newtheorem{lemma}[theorem]{Lemma}
\newtheorem{proposition}[theorem]{Proposition}
\newtheorem{corollary}[theorem]{Corollary}
\theoremstyle{definition}
\newtheorem{definition}[theorem]{Definition}
\theoremstyle{remark}
\newtheorem{remark}[theorem]{Remark}

\newcommand{\F}{\mathbb{F}}
\newcommand{\C}{\mathbb{C}}
\newcommand{\N}{\mathbb{N}}
\newcommand{\Z}{\mathbb{Z}}
\newcommand{\ID}{\mathrm{ID}}
\newcommand{\size}{\mathrm{size}^{\mathrm{nc}}}
\newcommand{\Hb}{\mathrm{H}}
\newcommand{\lean}[1]{\texttt{#1}}
\newcommand{\bits}{\F_2}
\DeclareMathOperator{\rank}{R}

\title{Majorana sums of squares and rectangular matrix multiplication:\\
noncommutative circuits for the identity polynomial of size $N^{1.6692}$}
\author{Produced by AI systems under the supervision of Carson Weitnauer\thanks{The mathematics, the Lean proofs
and this paper were produced by AI systems (OpenAI's Codex and Anthropic's Claude). Carson Weitnauer supervised the
research and is responsible for its release. How to attribute authorship of such work is an open question.}}
\date{October 29, 2026}

\begin{document}
\maketitle

\begin{abstract}
We prove that Hrube\v{s}'s identity polynomial $\ID_N=\sum_{i,j\in[N]}x_iy_jx_iy_j$ in $2N$ noncommuting
variables, the noncommutative elementary symmetric polynomial $S_{4,N}$, and the rectangular permanent
$\mathrm{perm}_{4,N}$ have noncommutative arithmetic circuits of size $O(N^{1.6692+\varepsilon})$ for every
$\varepsilon>0$, over every field of characteristic different from two. The previous published bound is
$O(N^{1.96})$, due to Hrube\v{s}. Without fast matrix multiplication, a different construction gives
$N^{1.72790+o(1)}$, also over every field of characteristic different from two.

The construction combines three ingredients: a weighted sum-of-squares identity over the bilinear forms
$X^{\mathsf T}P_\nu Y$ of real Pauli words $P_\nu$, whose weights are supported on a ``corner'' of the
Jordan--Wigner (Majorana) classes; blocks of Majorana modes that are moved by Clifford conjugation so that all forms
of a block are recovered by Walsh transforms of one rectangular matrix product; and a rectangular
matrix-multiplication bound. The
exponent is $2\Hb(\tfrac14)-2\mu\Hb(\tfrac1{4\mu})+\mu w$, where $w$ is an admissible exponent for multiplying
$n\times n^{(1-\mu)/\mu}$ by $n^{(1-\mu)/\mu}\times n$ matrices. With Uzna\'nski's 2026 bound
$\omega(1,k,1)\le1+k+\frac1{4k}$, in the integer form formalized by Zhou and Wang, this gives $1.66917$. With
OpenAI's 2026 bounds instead, an independent route gives $N^{1.690568+o(1)}$ in characteristic zero and in all but
finitely many characteristics.

The theorems and lemmas stated here are formalized and machine-checked in Lean~4 with Mathlib, using Zhou and
Wang's and OpenAI's Lean formalizations of the matrix-multiplication inputs and only Lean's standard axioms; the
remarks and numerical comparisons are not.
\end{abstract}

\section{Introduction}

\subsection{Background}
Hurwitz's problem~\cite{Hurwitz1898,Radon1922,Shapiro2000} asks for the smallest $s=\sigma(n)$ such that there is a
sum-of-squares identity
\begin{equation}\label{eq:sos}
(x_1^2+\dots+x_n^2)(y_1^2+\dots+y_n^2)=f_1^2+\dots+f_s^2
\end{equation}
with bilinear forms $f_k$. Hrube\v{s}, Wigderson and Yehudayoff~\cite{HWY2011} connected this classical question
to noncommutative arithmetic circuits, the model of Nisan~\cite{Nisan1991}: a lower bound $\sigma(n)\ge
n^{1+\varepsilon}$ would give an exponential lower bound for the noncommutative permanent. Their reduction runs
through the identity polynomial
\[
\ID_N=\sum_{i,j\in[N]}x_iy_jx_iy_j ,
\]
whose noncommutative circuit complexity is, up to constant factors, governed by sum-of-squares identities whose
bilinear forms are cheap to compute together (Hrube\v{s}~\cite[Lemma~22]{Hrubes2026}). In 2024 Hrube\v{s} proved the
first truly subquadratic bound $\sigma(n)=O(n^{1.62})$ over $\C$ and, combining it with rectangular matrix
multiplication, that $\ID_N$, $S_{4,N}$ and $\mathrm{perm}_{4,N}$ have noncommutative circuits of size
$O(N^{1.96})$ over every field of characteristic different from two~\cite{Hrubes2024,Hrubes2026}. (His proof gives
$O(N^{1.954})$ from the rectangular bound of Le~Gall and Urrutia~\cite{LeGallUrrutia2018}. Run with $n^r$ copies
of the $y$ variables, his argument gives the exponent $(r+\omega(1,r,1)+\delta)/(r+1)$ with $\delta<0.616$; with
Uzna\'nski's $\omega(1,2,1)\le\frac{25}{8}$~\cite{Uznanski2026}, the input behind Theorem~\ref{thm:main}, this is
below $1.9137$ at $r=2$. This comparison is computed here, not formalized.) The bound $\sigma(n)\le n^{\sqrt2+o(1)}$ over $\C$ follows from Appendix~B of Anthropic's October 2026
paper~\cite{Anthropic3SAT}; it was observed independently on October~27, 2026 by Kapila~\cite{Kapila2026}, in
this project (with a Lean proof built on Anthropic's formalization~\cite{AnthropicFormalMath}, in the
\texttt{theorem-b} directory of this paper's repository), and by
Tulsyan~\cite{Tulsyan2026}. Fed into Hrube\v{s}'s argument it improves the circuit exponent, but only to about
$1.85$--$1.89$ depending on the matrix-multiplication input. Before Hrube\v{s}, the best bound for $\ID_N$ was the
trivial $O(N^2)$, and branching programs of size about $N^2$ for $S_{4,N}$ and $\mathrm{perm}_{4,N}$ were given by
Arvind, Chatterjee, Datta and Mukhopadhyay~\cite{ACDM2019}.

\subsection{Results}
Throughout, $\size_\F(f)$ is the size of a smallest noncommutative circuit over the field $\F$ computing $f$
(Definition~\ref{def:circuit}), and we write $\size_\F(f_N)\le N^{c+o(1)}$ for
``$\size_\F(f_N)=O_\varepsilon(N^{c+\varepsilon})$ for every $\varepsilon>0$''.

\begin{theorem}[main theorem]\label{thm:main}
Over every field $\F$ of characteristic different from two, $\size_\F(\ID_N)$, $\size_\F(S_{4,N})$ and
$\size_\F(\mathrm{perm}_{4,N})$ are at most $N^{1.6692+o(1)}$.
\end{theorem}

The matrix-multiplication input is Uzna\'nski's rectangular bound~\cite{Uznanski2026}, through the integer rank
decompositions of Zhou and Wang~\cite{ZhouWang2026}: one integer decomposition maps into every field, so no
characteristic other than two is excluded. An independent route uses OpenAI's bounds instead.

\begin{theorem}[OpenAI's inputs]\label{thm:oairoute}
Let $c_0=211321/125000=1.690568$.
\begin{enumerate}[label=(\alph*)]
\item Over every field $\F$ of characteristic zero, the three polynomials have circuits of size $N^{c_0+o(1)}$.
\item There is a finite set $S$ of primes such that the same holds over every field $\F$ with
$\mathrm{char}\,\F\notin S\cup\{2\}$.
\end{enumerate}
\end{theorem}

The set $S$ in (b) comes from the Lefschetz principle and is not computed. Theorem~\ref{thm:oairoute} is weaker than
Theorem~\ref{thm:main}, but its matrix-multiplication input is independent of it, so either suffices for a
subquadratic bound below Theorem~\ref{thm:cubes}.

\begin{theorem}[no fast matrix multiplication]\label{thm:cubes}
Over every field of characteristic different from two, $\size(\ID_N)$, $\size(S_{4,N})$ and
$\size(\mathrm{perm}_{4,N})$ are at most $N^{1.727899+o(1)}$.
\end{theorem}
(The construction behind Theorem~\ref{thm:cubes} has exact exponent $c_\ast=1.72789810916\ldots$; the formal
statement uses the rational upper bound $1.727899$.) Its simplest precursor, a per-bit decomposition giving
$N^{\log_2(12)/2}=N^{1.7925}$, realizes the single-bit identity for the crossing tensor and the ``$\log 12$ of $4$
bits'' count in Appendix~H of~\cite{Anthropic3SAT}; Pauli (Bell) bases for the same tensor were also used by
Tulsyan~\cite{Tulsyan2026}, for a different cost measure.

Theorem~\ref{thm:cubes} is this project's earlier ``Pauli cube'' construction (October~28, 2026); we state it for
comparison and refer to its formalization (Section~\ref{sec:lean}). This paper proves Theorems~\ref{thm:main}
and~\ref{thm:oairoute}. Their heart is the following conditional bound, which turns block oracles for rectangular
products (Definition~\ref{def:oracle}) into a circuit exponent.

\begin{theorem}[the construction]\label{thm:construction}
Let $\F$ be a field with $2\ne0$ that contains a square root of $-1$. Let $\tfrac12<\mu<1$ and $w\ge2$, and put
$m=\lceil\mu n\rceil$. Suppose that for every $\varepsilon>0$ there is $C$ such that for every $n$ there is a block
oracle over $\F$ (Definition~\ref{def:oracle}) for $2^m\times2^{n-m}$ by $2^{n-m}\times2^m$ products with cost at
most $C(2^m)^{w+\varepsilon}$. Then
\[
\size_\F(\ID_N)\le N^{C(\mu,w)+o(1)},\qquad
C(\mu,w)=2\Hb\bigl(\tfrac14\bigr)-2\mu\,\Hb\Bigl(\frac1{4\mu}\Bigr)+\mu w ,
\]
where $\Hb(p)=-p\log_2p-(1-p)\log_2(1-p)$.
\end{theorem}

A block oracle is more than a matrix-multiplication algorithm over $\F$: its inputs are noncommuting, and it must
also evaluate a trilinear middle-insertion form. Two sources supply oracles of the required cost
(Section~\ref{sec:blocks}): straight-line programs over $\C$, so that over $\C$ any admissible exponent $w$ at aspect
$\frac{1-\mu}{\mu}$ in the sense of Section~\ref{sec:prelim-mm} suffices; and a single rank decomposition over $\F$
(Lemma~\ref{lem:scheme}), the source Theorem~\ref{thm:main} uses.

With Uzna\'nski's bound at $\mu=\frac{18}{29}$ (aspect $\frac{11}{18}$, $w=\frac{200}{99}$) one gets
$C=1.66916842\ldots$ (Section~\ref{sec:uz}). With OpenAI's bounds $\omega(1,0.465,1)=2$ and
$\omega(1,0.709,1)<2.092$~\cite{OAI2026b}, interpolated at $\mu=0.6337$, one gets $C<1.6905678<c_0$
(Section~\ref{sec:oai}); OpenAI's thin endpoint alone ($\mu=200/293$, $w=2$) gives $C=1.6938884\ldots$.

\begin{remark}
If $\omega(1,k,1)=2$ held for every $k\le1$, Theorem~\ref{thm:construction} would give exponents tending to
$2\Hb(\tfrac14)=1.62255\ldots$ as $\mu\to\tfrac12$, the growth exponent of the number of Pauli words in corner
classes (Section~\ref{sec:corner}). Since $\Hb\le1$, $C(\mu,w)\ge2\Hb(\frac14)$ whenever $w\ge2$, so this is the
floor of the present cost formula. It is not a lower bound on circuit size.
\end{remark}

\subsection{Techniques}
\emph{Corner sums of squares.} A weighted identity $\sum_\nu w_\nu f_\nu^2=(\sum X_x^2)(\sum Y_y^2)$ over the $4^n$
real Pauli bilinear forms $f_\nu(X,Y)=X^{\mathsf T}P_\nu Y$ holds if and only if a Fourier condition holds on the
isotropic labels (Lemma~\ref{lem:fourier}). In Jordan--Wigner coordinates~\cite{JordanWigner1928} a Pauli word is
a product of Majorana operators, described by a symmetric part $s$ and a skew part $a$. For weights depending only
on the class $(|s|,|a|)$ the Fourier condition becomes a small linear system with Krawtchouk
coefficients~\cite{Krawtchouk1929,MacWilliamsSloane1977}. We show it has a solution supported on the corner
$\min(a,n-a)+\min(b,n-b)\le n/2$ (Theorem~\ref{thm:corner}); the proof is a duality argument in which a dual
certificate would give a polynomial of degree at most $2n$ divisible by $(1-X^4)^{\lfloor n/2\rfloor+1}$.

\emph{Majorana blocks.} Swapping two Majorana modes is a conjugation by a single Clifford gate $1+cP_u$. A block of
$2m$ Majoranas can be moved onto the first $m$ qubits, after which every word supported in the block (up to a fixed
shift) has the form $P_{\alpha,\beta}\otimes I$. The corresponding bilinear forms are, after one shared change of
basis of the inputs, Walsh coefficients along shifted diagonals of a single $2^m\times 2^m$ matrix product of
$2^m\times 2^{n-m}$ and $2^{n-m}\times 2^m$ matrices. Hrube\v{s}'s middle-insertion term $f\star f^{yx}$ also
collapses to a single trilinear form. So one block costs one rectangular matrix product plus $O(m4^m)$.

\emph{Covering and entropy.} Each corner class is covered by blocks greedily, and an entropy calculation
(Lemma~\ref{lem:layout}) bounds the number of blocks by $2^{n(2\Hb(1/4)-2\mu\Hb(1/(4\mu)))+o(n)}$.

\emph{Matrix multiplication and characteristic.} Both routes reduce to one fixed rank decomposition of a matrix
tensor $\langle a,b,a\rangle$, from which a recursive bilinear algorithm supplies every block product. Zhou and
Wang's formalization provides integer decompositions along Uzna\'nski's curve, which work in every field. OpenAI's
bounds are statements about straight-line programs over $\C$: a program of cost $s$ gives a rank bound $2s$, ranks
multiply under tensor products, and one fixed product decomposition is an existential first-order statement with
integer coefficients, so it transfers by the Lefschetz principle.

\subsection{Formal verification}
The theorems above are formalized in Lean~4~\cite{Lean4} with Mathlib~\cite{Mathlib2020}. The formalization imports
OpenAI's Lean library for its matrix-multiplication bounds~\cite{OAI2026lean} unchanged, and Zhou and Wang's Lean
library~\cite{ZhouWang2026} with two declaration-level changes needed by our stricter compiler settings, so the
matrix-multiplication inputs are themselves machine-checked. The top-level Lean theorems, the definitions they use, and the
checks performed are listed in Section~\ref{sec:lean}.

\section{Preliminaries}

\subsection{Noncommutative circuits}
\begin{definition}[Hrube\v{s}~\cite{Hrubes2026}]\label{def:circuit}
A \emph{noncommutative circuit} over $\F$ is a directed acyclic graph whose in-degree-zero nodes are labelled by
variables or elements of $\F$, and whose other nodes have in-degree two and are labelled $+$ or $\times$, the two
edges into a $\times$ node being labelled left and right. Each node computes a polynomial in the free algebra
$\F\langle x_1,\dots\rangle$; the circuit computes $f$ if some node does; its size is its number of nodes.
\end{definition}

The polynomials of interest are $\ID_N$ above,
\[
S_{4,N}=\sum_{(i_1,\dots,i_4)}x_{i_1}x_{i_2}x_{i_3}x_{i_4},\qquad
\mathrm{perm}_{4,N}=\sum_{(i_1,\dots,i_4)}x_{1,i_1}x_{2,i_2}x_{3,i_3}x_{4,i_4},
\]
with sums over ordered quadruples of pairwise distinct elements of $[N]$. Hrube\v{s}~\cite[Cor.~27]{Hrubes2026}
reduces $S_{4,N}$ and $\mathrm{perm}_{4,N}$ to $\ID_N$ with linear overhead (through $\ID^\ast$ and
Proposition~26 there); we use these reductions as they are, so it suffices to bound $\size(\ID_N)$. Since setting
variables to zero does not increase size, it also suffices to treat $N=2^n$, indexing the variables by bit vectors
$x\in\F_2^n$.

\subsection{From sums of squares to circuits}
For a bilinear form $f=\sum_{x,y}a_{xy}X_xY_y$ and a new variable $z$, Hrube\v{s}~\cite{Hrubes2026} defines the
noncommutative polynomials
\[
f^{xy}=\sum a_{xy}X_xY_y,\quad f^{yx}=\sum a_{xy}Y_yX_x,\quad f\star z=\sum a_{xy}X_xzY_y,\quad
f^{[2]}=\tfrac12\bigl(f^{xy}f^{xy}+f\star f^{yx}\bigr).
\]
\begin{lemma}[Hrube\v{s}, proof of~{\cite[Lemma~22]{Hrubes2026}}]\label{lem:hrubes}
Let $2\ne0$ in $\F$. If $(\sum_xX_x^2)(\sum_yY_y^2)=\sum_kw_kf_k^2$ as commutative polynomials, with bilinear $f_k$,
then $\sum_kw_kf_k^{[2]}=\ID$ in the free algebra.
\end{lemma}
The cost of $\sum_k w_kf_k^{[2]}$ is the cost of the forms $f^{xy}_k$ and $f^{yx}_k$ plus that of the single
trilinear form $\sum_kw_kf_k\star z_k$ evaluated at $z_k=f_k^{yx}$; Hrube\v{s} bounds the latter by a
Baur--Strassen argument~\cite{BaurStrassen1983}. Our blocks make all three cheap at once (Section~\ref{sec:blocks}).

\subsection{Pauli forms in bit coordinates}
A \emph{word} is $\nu\in\{I,X,Z,J\}^n$, where $J=ZX=-XZ$; write $x_\nu,z_\nu\in\F_2^n$ for its $X$- and $Z$-parts. Its
real Pauli matrix, indexed by $\F_2^n$, is
\[
P_\nu[x,y]=[\,y=x+x_\nu\,]\,(-1)^{z_\nu\cdot x},
\]
and we write $P_{\alpha,\beta}$ for the matrix with $X$-part $\alpha$ and $Z$-part $\beta$. These satisfy
$P_{\alpha,\beta}P_{\alpha',\beta'}=(-1)^{\beta'\cdot\alpha}P_{\alpha+\alpha',\beta+\beta'}$. The \emph{Pauli
form} of $\nu$ is $f_\nu(X,Y)=X^{\mathsf T}P_\nu Y$.

\subsection{Rectangular matrix multiplication}\label{sec:prelim-mm}
For $0\le k\le1$, an exponent $\tau$ is \emph{admissible at aspect $k$} if for every $\varepsilon>0$ there is $C$
such that for every $n\ge1$ some straight-line program (additions, subtractions, multiplications, constants)
multiplies $n\times\lceil n^k\rceil$ by $\lceil n^k\rceil\times n$ matrices with at most $C n^{\tau+\varepsilon}$
operations. We write $\omega(1,k,1)$ for the infimum of admissible exponents and $\alpha=\sup\{k:\omega(1,k,1)=2\}$.
These are OpenAI's definitions~\cite{OAI2026lean}, which our Lean development imports.

\section{The corner sum of squares}\label{sec:corner}

\begin{lemma}[Fourier criterion]\label{lem:fourier}
Let $2\ne0$ in $\F$ and $w:\{\text{words}\}\to\F$. Put
\[\hat w(t,b)=\sum_\nu w_\nu(-1)^{z_\nu\cdot t+x_\nu\cdot b}.\]
If $\hat w(t,b)=2^n[t=b=0]$ for all $t,b\in\F_2^n$ with $t\cdot b=0$, then
$(\sum_xX_x^2)(\sum_yY_y^2)=\sum_\nu w_\nu f_\nu(X,Y)^2$.
\end{lemma}
\begin{proof}
Comparing coefficients of $X_xY_yX_uY_v$, the identity is equivalent to
$\sum_\nu w_\nu(P_\nu[x,y]P_\nu[u,v]+P_\nu[x,v]P_\nu[u,y])=2[x=u,\,y=v]$. Both products vanish unless
$d:=x+u=y+v$, and then the left side is $g(x+y,d)+g(x+y+d,d)$ with
$g(\alpha,d)=\sum_{x_\nu=\alpha}w_\nu(-1)^{z_\nu\cdot d}$. Taking the Fourier transform in $\alpha$, the required
identity $g(\alpha,d)+g(\alpha+d,d)=2[d=0]$ reads $\hat w(d,b)(1+(-1)^{b\cdot d})=2\cdot2^n[d=0][b=0]$, which is
automatic when $b\cdot d=1$ and is the hypothesis when $b\cdot d=0$.
\end{proof}

\subsection{Majorana coordinates}
Jordan--Wigner coordinates~\cite{JordanWigner1928} write each word as $\nu=\nu(s,a)$ with $s,a\in\F_2^n$ via
\[
x_\nu=s+a,\qquad (z_\nu)_j=a_j+\sum_{k>j}(s+a)_k .
\]
This is a bijection; $P_{\nu(s,a)}$ is, up to sign, the product of the symmetric Majoranas $Z\cdots ZX_j$ over the
support of $s$ and the skew Majoranas $Z\cdots ZJ_j$ over the support of $a$. The \emph{class} of $\nu$ is
$(|s|,|a|)$. In these coordinates the commutation form is
$\langle(s,a),(s',a')\rangle=(|s|+|a|)(|s'|+|a'|)+s\cdot s'+a\cdot a'$, and the label $(t,b)$ with Majorana
coordinates $(s',a')$ is isotropic ($t\cdot b=0$) if and only if
$q(|s'|,|a'|):=|a'|+\binom{|s'|+|a'|}{2}\equiv0\pmod 2$.

For class-constant weights $w_\nu=2^nW(|s|,|a|)$, summing characters over a class gives Krawtchouk
numbers $K_a(c)=[z^a](1-z)^c(1+z)^{n-c}$, and the Fourier criterion becomes the finite linear system
\begin{equation}\label{eq:system}
\sum_{a,b=0}^nW(a,b)\,(-1)^{(a+b)(c+d)}K_a(c)K_b(d)=[c=d=0]\qquad(0\le c,d\le n,\ q(c,d)=0).
\end{equation}

\begin{theorem}[corner feasibility]\label{thm:corner}
Let $\F$ have $2\ne0$ and contain $i$ with $i^2=-1$. For every $n$, the system \eqref{eq:system} has a solution
$W$ supported on the \emph{corner}
\[
\mathcal C_n=\{(a,b):\ \min(a,n-a)+\min(b,n-b)\le\lfloor n/2\rfloor\}.
\]
Consequently there are weights on the $4^n$ words, supported on corner classes, with
$(\sum X_x^2)(\sum Y_y^2)=\sum_\nu w_\nu f_\nu(X,Y)^2$.
\end{theorem}
\begin{proof}
If not, linear duality gives $y(c,d)$, supported on the pairs with $q(c,d)=0$, with $y(0,0)\ne0$ and
$\sum_{c,d}y(c,d)(-1)^{(a+b)(c+d)}K_a(c)K_b(d)=0$ for every $(a,b)\in\mathcal C_n$. Put
$P_a(X)=(1-X)^a(1+X)^{n-a}+(1+X)^a(1-X)^{n-a}$ and
$\Psi=\sum_{a,b}\mathrm{coef}(a,b)\,P_aP_b$, where $\mathrm{coef}(a,b)$ is the left side above.

\emph{$\Psi$ is divisible by $(1-X^2)^m$, $m=\lfloor n/2\rfloor+1$.} For $(a,b)\in\mathcal C_n$ the coefficient
vanishes; otherwise $((1-X)(1+X))^{\min(a,n-a)+\min(b,n-b)}$ divides $P_aP_b$, and the exponent is at least $m$.

\emph{$\Psi$ is a multiple of $G=\sum_{c,d\ \text{even}}y(c,d)X^{c+d}$.} Exchanging the sums gives
$\Psi=\sum_{c,d}y(c,d)\,S_{c,c+d}S_{d,c+d}$, where $S_{c,e}=\sum_a(-1)^{ae}K_a(c)P_a$. The homogenization
identity $\sum_aK_a(c)s^at^{n-a}=(t-s)^c(t+s)^{n-c}$ gives
\[
S_{c,e}=\begin{cases}\bigl((2X)^c+(-2X)^c\bigr)\,2^{n-c}&e\text{ even},\\
2^c\bigl((2X)^{n-c}+(-2X)^{n-c}\bigr)&e\text{ odd},\end{cases}
\]
so $S_{c,c+d}S_{d,c+d}=4\cdot2^{2n}X^{c+d}$ when $c$ and $d$ are
both even and $0$ otherwise. Hence $\Psi=4\cdot2^{2n}G$.

\emph{$G$ is divisible by $(1+X^2)^m$.} For even $c,d$, $q(c,d)=0$ forces $4\mid c+d$, so $G(iX)=G(X)$; applying
$X\mapsto iX$ to $(1-X^2)^m\mid G$ gives $(1+X^2)^m\mid G$.

Since $2\ne0$, $(1-X^2)^m$ and $(1+X^2)^m$ are coprime, so $(1-X^4)^m\mid G$. But $\deg G\le2n<4m$, so $G=0$
and $y(0,0)=G(0)=0$, a contradiction. The second sentence follows from Lemma~\ref{lem:fourier}.
\end{proof}

The number of words in corner classes is $\sum_{(a,b)\in\mathcal C_n}\binom na\binom nb=2^{2\Hb(1/4)n+O(\log n)}$:
the maximum is at $a,b\approx n/4$.

\section{Majorana blocks}\label{sec:blocks}

\subsection{Moving Majoranas by conjugation}
Index the $2n$ Majoranas by $k\in\{1,\dots,2n\}$ and identify a word with its Majorana vector $A\in\F_2^{2n}$ (the
pair $(s,a)$); write $P_A$ for its matrix. For $k\ne l$ let $u=e_k+e_l$, choose $c\in\{1,i\}$ with
$(cP_u)^2=-1$, and put $G_{kl}=1+cP_u$, so $G_{kl}^{-1}=\tfrac12(1-cP_u)$.

\begin{lemma}\label{lem:swap}
$G_{kl}P_AG_{kl}^{-1}=\lambda\,P_{\mathrm{swap}_{kl}(A)}$ for some $\lambda\ne0$.
\end{lemma}
\begin{proof}
$P_u$ and $P_A$ commute or anticommute according to $\langle u,A\rangle=A_k+A_l$. If they commute, conjugation
fixes $P_A$ and $\mathrm{swap}_{kl}$ fixes $A$. If they anticommute, $G P_AG^{-1}=-P_A\,cP_u=\lambda P_{A+u}$, and
$A+u=\mathrm{swap}_{kl}(A)$.
\end{proof}

A \emph{block} is a set $U$ of $2m$ Majoranas; the \emph{canonical block} $T$ consists of the $2m$ Majoranas of
the first $m$ qubits. Composing at most $|U\setminus T|\le2m$ swaps gives a matrix $\Gamma$ with
$\Gamma P_A\Gamma^{-1}=\lambda_AP_{A'}$, $A'$ supported in $T$, for every $A$ supported in $U$. For a fixed
\emph{shift} $\sigma\in\F_2^{2n}$, $P_{A+\sigma}=\pm P_AP_\sigma$, hence
\begin{equation}\label{eq:blockword}
P_{A+\sigma}=\lambda\,\Gamma^{-1}P_{A'}\,(\Gamma P_\sigma)\qquad(\lambda\ne0).
\end{equation}

\subsection{One block is one rectangular product}
Since $\Gamma^{-1}$ and $\Gamma P_\sigma$ are matrices with entries in the coefficient field,
\eqref{eq:blockword} transforms all three
noncommutative forms of a block word at once: with $X'=(\Gamma^{-1})^{\mathsf T}X$ and $Y'=\Gamma P_\sigma Y$,
\[
f^{xy}_{A+\sigma}=\lambda\,f^{xy}_{A'}(X',Y'),\quad f^{yx}_{A+\sigma}=\lambda\,f^{yx}_{A'}(X',Y'),\quad
f_{A+\sigma}\star z=\lambda\,(f_{A'}\star z)(X',Y').
\]
Each gate factor $c_1+c_2P$ has two nonzero entries per row, so $X'$ and $Y'$ cost $O(m2^n)$ gates. Since $A'$ is
supported on the first $m$ qubits, $P_{A'}=P_{\alpha,\beta}$ with $\alpha,\beta$ supported there; splitting
$x=(p,s)\in\F_2^m\times\F_2^{n-m}$,
\[
f^{xy}_{\alpha,\beta}(X',Y')=\sum_{p}(-1)^{\beta\cdot p}M[p,p+\alpha],\qquad
M[p,q]=\sum_{s\in\F_2^{n-m}}X'(p,s)\,Y'(q,s).
\]
So the $xy$-forms of all words of the block are Walsh coefficients along shifted diagonals of one product $M$ of a
$2^m\times2^{n-m}$ and a $2^{n-m}\times2^m$ matrix; the $yx$-forms come from the swapped product $M'$. Grouping the
block's words by their canonical label $(\alpha,\beta)$, the middle-insertion terms combine into one trilinear
form $\sum_{p,s,q}X'(p,s)\,Q[p,q]\,Y'(q,s)$ with $Q[p,q]=\sum_b(-1)^{p\cdot b}g(p+q,b)$ for suitable values $g$.

\begin{definition}\label{def:oracle}
A \emph{block oracle} with cost $M$ for $2^m\times2^r$ by $2^r\times2^m$ products is a family of circuits that, from
arbitrary available values $X'(p,s)$, $Y'(q,s)$ and $Q[p,q]$ in the free algebra ($p,q\in\F_2^m$, $s\in\F_2^r$),
compute all entries $\sum_sX'(p,s)Y'(q,s)$, and the single value $\sum_{p,s,q}X'(p,s)Q[p,q]Y'(q,s)$, each with at
most $M$ further gates.
\end{definition}

\begin{proposition}\label{prop:block}
Given a block oracle with cost $M$, the sum $\sum_\nu c_\nu f_\nu^{[2]}$ over all words of a block, with arbitrary
coefficients, takes at most $O(m2^n)+3M+15(m+1)4^m+2$ gates, independently of the number of words.
\end{proposition}
\begin{proof}
Products and swapped products: $2M$. Walsh--Hadamard transforms of all $2^m$ shifted diagonals: $O(m4^m)$ via the
Kronecker structure. Squares and the regrouped coefficients: $O(4^m)$. Middle insertion: $M$, after $O(m4^m)$ to
form $Q$.
\end{proof}

Block oracles come from two sources.
\begin{enumerate}[label=(\roman*)]
\item \emph{Programs (over $\C$).} A correct straight-line program of cost $s$ for the
$2^m\times2^{n-m}\times2^m$ product gives a block oracle with cost $960s$. Each register is split into its constant,
$x$-linear, $y$-linear and \emph{ordered} $xy$-bilinear parts (the product rule
$\partial^2(pq)=p(0)\partial^2q+q(0)\partial^2p+\partial_xp\,\partial_yq+\partial_xq\,\partial_yp$ keeps the $x$
factor on the left), which yields all ordered products with $91s$ gates; the middle insertion is obtained by reverse
differentiation of the program, in the spirit of Baur--Strassen~\cite{BaurStrassen1983} and Hrube\v{s}'s
Lemma~20~\cite{Hrubes2026}. The polynomial semantics uses that $\C$ is infinite.
\item \emph{Rank decompositions (any field).} A decomposition
$\langle a,b,a\rangle=\sum_{r\le R}u_r\otimes v_r\otimes w_r$ gives the classical recursive bilinear algorithm,
which is valid in any noncommutative algebra because scalars are central and $x$-factors stay on the left; the
trilinear identity
$\sum_{i,j,l}X_{ij}Q_{il}Y_{jl}=\sum_r\bigl(\sum u_rX\bigr)\bigl(\sum w_rQ\bigr)\bigl(\sum v_rY\bigr)$ recurses the
same way. With $t$ levels this costs at most $(7(2ab+a^2)+6)(t+1)\max(R,ab,a^2)^t$.
\end{enumerate}

\section{Covering the corner classes}\label{sec:cover}

Fix $\tfrac12<\mu<1$ and $m=\lceil\mu n\rceil$. The words of class $(a,b)$ are exactly those with Majorana vector
$A+\sigma_{a,b}$, where $\sigma_{a,b}$ complements the symmetric part if $a>n-a$ and the skew part if $b>n-b$, and
$A$ has weights $(a_0,b_0)=(\min(a,n-a),\min(b,n-b))$. A block with $u_s$ symmetric and $u_a$ skew Majoranas
($u_s+u_a=2m$) contains $\binom{u_s}{a_0}\binom{u_a}{b_0}$ of them.

\begin{lemma}[greedy cover]\label{lem:cover}
For $a_0\le u_s\le n$ and $b_0\le u_a\le n$ there is a family of at most
\[
1+2n\,\frac{\binom n{a_0}\binom n{b_0}}{\binom{u_s}{a_0}\binom{u_a}{b_0}}
\]
blocks $(U_s,U_a)$ with $|U_s|=u_s$, $|U_a|=u_a$ such that every pair $(S,T)$ with $|S|=a_0$, $|T|=b_0$ has
$S\subseteq U_s$, $T\subseteq U_a$ for some block of the family.
\end{lemma}
\begin{proof}
Every pair lies in $D=\binom{n-a_0}{u_s-a_0}\binom{n-b_0}{u_a-b_0}$ of the $\binom n{u_s}\binom n{u_a}$ blocks, and
$\binom n{u}\binom u{a_0}=\binom n{a_0}\binom{n-a_0}{u-a_0}$. The greedy algorithm covers at least a fraction
$D/\#\text{blocks}$ of the remaining pairs at each step, so $1+(\#\text{blocks}/D)\ln\#\text{pairs}$ blocks
suffice, and $\ln\#\text{pairs}\le 2n$.
\end{proof}

\begin{lemma}[layout]\label{lem:layout}
Let $g_\mu=2\Hb(\tfrac14)-2\mu\Hb(\tfrac1{4\mu})$. For every $\delta>0$ and all large $n$, every corner class
$(a,b)\in\mathcal C_n$ admits $u_s+u_a=2m$, $u_s,u_a\le n$, $a_0\le u_s$, $b_0\le u_a$, with
\[
\binom na\binom nb\le2^{n(g_\mu+\delta)}\binom{u_s}{a_0}\binom{u_a}{b_0}.
\]
\end{lemma}
\begin{proof}[Proof sketch]
Write $L=2\mu$, $\alpha=a_0/n\le\beta=b_0/n$, $t=\alpha+\beta\le\frac12$. If $\beta/t\le1/L$, take
$u_s/n=L\alpha/t$, $u_a/n=L\beta/t$; concavity of entropy bounds the ratio by
$2^{n(2\Hb(t/2)-L\Hb(t/L))}$, and $2\Hb(t/2)-L\Hb(t/L)$ is nondecreasing in $t$, so it is at most $g_\mu$. If
$\beta/t>1/L$, take $u_a=n$, $u_s=(L-1)n$; the exponent $\Hb(\alpha)-(L-1)\Hb(\alpha/(L-1))$ is increasing in $\alpha$
and agrees with the first case on the boundary. The formal proof works with exact binomials: a Vandermonde-type
identity $\sum_u\binom u{a_0}\binom{M-u}{b_0}=\binom{M+1}{a_0+b_0+1}$ chooses the integer layout, the bounds
$\binom NK K^K(N-K)^{N-K}\le N^N\le(N+1)\binom NKK^K(N-K)^{N-K}$ convert to entropy with only a polynomial loss,
and convexity of $x\log x$ reduces the exponent to the extreme case $a_0+b_0=\lfloor n/2\rfloor$; the threshold
is $n_0=\lceil54/(\delta\ln2)^3\rceil+1$.
\end{proof}

\section{Assembling the circuit}\label{sec:assembly}

\begin{proof}[Proof of Theorem~\ref{thm:construction}]
Let $N=2^n$. By Theorem~\ref{thm:corner} and Lemma~\ref{lem:hrubes}, $\ID=\sum_\nu w_\nu f_\nu^{[2]}$ with $w$
supported on corner classes. For each corner class choose the layout of Lemma~\ref{lem:layout} and the cover of
Lemma~\ref{lem:cover}; assign each word to one covering block. By Proposition~\ref{prop:block} each block costs
$O(m2^n)+3M+O(m4^m)$, with $M\le C_\varepsilon(2^m)^{w+\varepsilon}$ by hypothesis. There are at most $(n+1)^2$
classes and, per class, at most $1+2n2^{n(g_\mu+\delta)}$ blocks. Since $\mu>\frac12$ and $w\ge2$,
$2^n\le4^m\le(2^m)^w$, so the total is
\[
2^{n(g_\mu+\delta)+\mu n(w+\varepsilon)+O(\log n)}=N^{C(\mu,w)+o(1)}.
\]
Arbitrary $N$ follows by padding to the next power of two.
\end{proof}

\section{Matrix multiplication inputs}\label{sec:mm}

Both routes supply the block oracles of Theorem~\ref{thm:construction} from one fixed rank decomposition, through
source~(ii) of Section~\ref{sec:blocks}, and both pass to the base field through the regular representation.

\begin{lemma}[one decomposition suffices]\label{lem:scheme}
Let $\frac12<\mu<1$ and $k=\frac{1-\mu}{\mu}$. Let $E$ be a field with $2\ne0$ that contains a square root of $-1$.
Suppose that over $E$ the tensor $\langle a,b,a\rangle$, with $a\ge2$ and $b\ge a^k$, has rank at most $\rho$,
and that $\max(\rho,ab,a^2)\le a^\tau$. Then $\size_E(\ID_N)\le N^{C(\mu,\tau)+o(1)}$.
\end{lemma}
\begin{proof}
Let $m=\lceil\mu n\rceil$ and let $t$ be least with $a^t\ge2^m$. Since $n-m\le mk$, also
$b^t\ge a^{kt}\ge2^{n-m}$, so the block products of size $2^m\times2^{n-m}\times2^m$ embed, by padding with zeros,
in $\langle a^t,b^t,a^t\rangle$. Source~(ii) gives block oracles with cost at most
$(7(2ab+a^2)+6)(t+1)\,a^{t\tau}$, and $a^{t\tau}\le a^\tau(2^m)^\tau$. Since $t=O(m)$, this is
$O((2^m)^{\tau+\varepsilon})$ for every $\varepsilon>0$, and Theorem~\ref{thm:construction} applies.
\end{proof}

\begin{lemma}[descent]\label{lem:descent}
If $E/\F$ is a finite extension of degree $d$, then $\size_\F(\ID_N)\le d^2(2d+1)\,\size_E(\ID_N)$.
\end{lemma}
\begin{proof}
$E$ embeds in $d\times d$ matrices over $\F$ by its regular representation, so a circuit over $E$ becomes a
circuit over $\F$ computing a $d\times d$ matrix of noncommutative polynomials, each gate simulated entrywise with at
most $d^2(2d+1)$ gates. A diagonal entry of the image of $\ID$ is $\ID$ itself.
\end{proof}

\subsection{Uzna\'nski's curve over the integers}\label{sec:uz}
Let $f(k)=2$ for $0\le k\le\frac12$ and $f(k)=1+k+\frac1{4k}$ for $k>\frac12$. Uzna\'nski~\cite{Uznanski2026}
proved $\omega(1,k,1)\le f(k)$; in particular $\alpha\ge\frac12$ and $\omega\le\frac94$, the bound OpenAI had proved
over $\C$ shortly before~\cite{OAI2026a}. Zhou and Wang~\cite{ZhouWang2026} extended the bound to every associative
ring, with integer bilinear algorithms, and formalized the extension in Lean. We use one theorem of their
formalization.

\begin{theorem}[Uzna\'nski~\cite{Uznanski2026}; Zhou and Wang~\cite{ZhouWang2026}]\label{thm:uz}
For every $k\ge0$ and $\varepsilon>0$ there are integers $a\ge2$, $b\ge a^k$ and $\rho\ge\max(ab,a^2)$, and
integer vectors $u_r,v_r,w_r$ $(r\le\rho)$, such that $\langle a,b,a\rangle=\sum_{r\le\rho}u_r\otimes v_r\otimes
w_r$ over $\Z$ and $\log\rho/\log a<f(k)+\varepsilon$.
\end{theorem}

In Lean this is \lean{Matmul.RectangularRankBridge.exists\_scheme\_certificate}. Its proof bounds the asymptotic
rank of integer matrix tensors through spectral points in the sense of Strassen~\cite{Strassen1988}, following
Uzna\'nski's shared-leg entropy argument, and extracts a single decomposition with the given slack.

\begin{proof}[Proof of Theorem~\ref{thm:main}]
Take $\mu=\frac{18}{29}$. Then $k=\frac{1-\mu}{\mu}=\frac{11}{18}$ and $f(k)=1+\frac{11}{18}+\frac{9}{22}=\frac{200}{99}$.
Let $\varepsilon_0=10^{-5}$ and fix the decomposition of Theorem~\ref{thm:uz} for this $k$ and $\varepsilon=\varepsilon_0$.

Let $\F$ have characteristic other than two, and let $E$ be $\F$ with a root of $X^2+1$ adjoined (or $E=\F$ if
$\F$ already has one), so $[E:\F]\le2$. The integer decomposition maps into $E$, and
$\max(\rho,ab,a^2)=\rho\le a^{\tau}$ with $\tau=\frac{200}{99}+\varepsilon_0$. Lemma~\ref{lem:scheme} gives
$\size_E(\ID_N)\le N^{C(\mu,\tau)+o(1)}$, where, since $C$ is affine in $w$,
\[
C(\mu,\tau)=C\bigl(\tfrac{18}{29},\tfrac{200}{99}\bigr)+\tfrac{18}{29}\varepsilon_0
=1.66916842\ldots+0.0000062\ldots<1.6692 .
\]
The value $C(\frac{18}{29},\frac{200}{99})+10^{-5}<1.6692$ is certified with rational enclosures of $\log2$,
$\log\frac32$, $\log\frac{36}{29}$ and $\log\frac{72}{43}$. Lemma~\ref{lem:descent} transfers the bound to $\F$, and
Hrube\v{s}'s reductions give $S_{4,N}$ and $\mathrm{perm}_{4,N}$.
\end{proof}

The choice $\mu=\frac{18}{29}$ is a rational point near the optimum: over $\mu$, the exponent with $w=f(k)$ is
minimized near $\mu\approx0.6207$, at $1.669168\ldots$.

\subsection{OpenAI's bounds}\label{sec:oai}

\begin{theorem}[OpenAI~\cite{OAI2026b,OAI2026lean}]\label{thm:oai}
Over $\C$: $\alpha>0.465$, so $\omega(1,\tfrac{93}{200},1)=2$; and $\omega(1,0.709,1)<\tfrac{523}{250}=2.092$.
\end{theorem}
These are formalized in OpenAI's Lean library, which we import unchanged. The first is used through the strict
supremum: some aspect $k>0.465$ has $\omega(1,k,1)=2$, and admissibility is monotone in the aspect.

\begin{lemma}[programs give ranks]\label{lem:rank}
A correct straight-line program of cost $s$ for $a\times b$ by $b\times c$ matrix multiplication over $\C$ yields
a rank decomposition of the matrix-multiplication tensor $\langle a,b,c\rangle$ of size $2s$.
\end{lemma}
\begin{proof}
This is the classical comparison of multiplicative and bilinear complexity
(Strassen~\cite{Strassen1973}; see~\cite[Ch.~14]{BCS1997}): by the product rule above, the bilinear part of every
register is a linear combination of the $2\cdot\#\{\times\text{ gates}\}$ rank-one forms
$\partial_xp\otimes\partial_yq$, $\partial_xq\otimes\partial_yp$, and correctness identifies the bilinear parts of
the outputs with the tensor.
\end{proof}

\begin{lemma}[rational interpolation]\label{lem:interp}
If $\tau_i\ge2$ is admissible at aspect $k_i\in[0,1]$ $(i=1,2)$ over $\C$, then for $p,q\in\N$ with $p+q>0$ the
exponent $\frac{p\tau_1+q\tau_2}{p+q}$ is admissible at aspect $\frac{pk_1+qk_2}{p+q}$.
\end{lemma}
This is a special case of the classical convexity of $\omega(1,k,1)$ (Lotti and Romani~\cite{LottiRomani1983}).
\begin{proof}
Fix $\varepsilon>0$. For large $t$ take programs for sizes $n_1=2^{tp}$ and $n_2=2^{tq}$ at the two endpoints,
convert them to rank decompositions (Lemma~\ref{lem:rank}), and take their tensor product: a decomposition of
$\langle a,b,a\rangle$ with $a=2^{t(p+q)}$, $b\ge a^{k}$, and rank at most
$4C_1C_2\,a^{\tau+\varepsilon}\le a^{\tau+2\varepsilon}$ once $t$ is large. The block recursion
($\langle a,b,a\rangle\otimes$ a program for any size) turns this single decomposition into an admissible exponent
$\tau+2\varepsilon$.
\end{proof}

With $p=829933$ and $q=716295$, the endpoints $(\frac{93}{200},2)$ and $(\frac{709}{1000},\frac{523}{250})$
interpolate to aspect $k^\ast=\frac{3663}{6337}=\frac{1-\mu}{\mu}$ for $\mu=\frac{6337}{10000}$, with exponent
$w^\ast=\frac{157917757}{77311400}=2.0426192\ldots$. Then
\[
C(\mu,w^\ast)=1.69056770515\ldots<1.690568-10^{-7},
\]
certified with rational enclosures of $\log2$, $\log\frac32$, $\log\frac{10000}{6337}$ and
$\log\frac{7674}{6337}$. This proves Theorem~\ref{thm:oairoute}(a) over $\C$; every field of characteristic zero is
covered by (b), whose set $S$ contains only primes.

\begin{remark}
The same formula with only the rectangular bounds of Alman et al.~\cite{ADVWXXZ2024} gives $C<1.7118$, and
OpenAI's thin endpoint alone ($\mu=\frac{200}{293}$, $w=2$) gives $1.6938884\ldots$; the latter is also formalized.
\end{remark}

\subsection{Other characteristics for OpenAI's route}\label{sec:charp}

\begin{proof}[Proof of Theorem~\ref{thm:oairoute}(b)]
Let $\varepsilon_0=10^{-8}$. The proof of Lemma~\ref{lem:interp} produces one rank decomposition over $\C$ of a
fixed tensor $\langle a,b,a\rangle$, with $a^{k^\ast}\le b$ and $\max(R,ab,a^2)\le a^{w^\ast+2\varepsilon_0}$.
Its existence is an existential first-order sentence in the language of rings with integer coefficients, true in
$\C$. By completeness of the theory of algebraically closed fields of characteristic zero and the Lefschetz
principle, it holds in every algebraically closed field of characteristic $0$ or $p\notin S$, for a finite set $S$
of primes.

Let $\F$ have characteristic $0$ or $p\notin S\cup\{2\}$. The decomposition exists over the algebraic closure; its
finitely many coefficients and a square root of $-1$ generate a finite extension $E/\F$. Lemma~\ref{lem:scheme}
gives $\size_E(\ID_N)\le N^{C(\mu,w^\ast+2\varepsilon_0)+o(1)}$, and
$C(\mu,w^\ast+2\varepsilon_0)=C(\mu,w^\ast)+2\mu\varepsilon_0<c_0$. Lemma~\ref{lem:descent} transfers the bound to
$\F$, and Hrube\v{s}'s reductions give $S_{4,N}$ and $\mathrm{perm}_{4,N}$.
\end{proof}

The argument is not specific to the numerical choice.

\begin{theorem}[parametric field transfer]\label{thm:fieldtransfer}
Let $\frac12<\mu<1$ and let $w$ be an admissible rectangular exponent over $\C$ at aspect $\frac{1-\mu}{\mu}$.
\begin{enumerate}[label=(\alph*)]
\item If $c\ge1$ and $C(\mu,w)<c$, there is a finite set $S$ of primes such that the three circuit bounds hold at
exponent $c$ over every field of characteristic zero or characteristic outside $S\cup\{2\}$. The set $S$ does not
depend on the field or on the slack $\varepsilon$ in the circuit bound.
\item If $C(\mu,w)\ge1$, the three bounds hold at the exact exponent $C(\mu,w)$ over every field of characteristic
zero, whether or not it contains a square root of $-1$.
\end{enumerate}
\end{theorem}
\begin{proof}
For (a), put $\eta=(c-C(\mu,w))/(4\mu)>0$. The proof of Lemma~\ref{lem:interp}, with one endpoint, gives one rank
decomposition with exponent at most $w+2\eta$; transfer it and descend as above. The resulting exponent is
$C(\mu,w+2\eta)=C(\mu,w)+2\mu\eta=\frac12(C(\mu,w)+c)<c$. For (b), apply (a) with $c=C(\mu,w)+\delta$ for every
$\delta>0$; a field of characteristic zero avoids every exceptional set. Given slack $\varepsilon$, use
$\delta=\varepsilon/2$ and slack $\varepsilon/2$.
\end{proof}
Part~(b) does not give one finite exceptional set in positive characteristic at the exact exponent: the sets
produced as $\delta\to0$ need not have a finite union.

No transfer of $\alpha>0.465$ itself to positive characteristic is known; OpenAI states explicitly that none is
asserted~\cite[\S14]{OAI2026b}, because $\alpha$ requires decompositions of every accuracy. This argument needs only
one, because the exponent has slack. This is also why $S$ is not computed. Theorem~\ref{thm:main} makes the
question moot for the circuit bound, since its decompositions have integer coefficients from the start.

\section{Formal verification}\label{sec:lean}

\subsection{What is checked}
The development uses Lean~4.34.1 and Mathlib (commit \texttt{d13f23b}). It consists of
\begin{itemize}
\item OpenAI's matrix-multiplication library~\cite{OAI2026lean} (509 files from \texttt{openai/math}, commit
\texttt{7ef5c0c}), unchanged, with their SHA-256 hashes recorded;
\item the \lean{Matmul} library of Zhou and Wang~\cite{ZhouWang2026} (64 files, commit \texttt{285f513}), with its
license and notice files. Two of its declarations relied on Lean's automatic implicit arguments, which our build
disables; we added the missing binder \lean{\{t : Nat\}} to each, checked that the resulting types and bodies are
identical to the upstream elaboration, and marked the file as changed. Nothing else was changed;
\item harfe's formalization of Brouwer's fixed-point theorem~\cite{harfe2026}, a dependency of both libraries,
adapted to our Lean and Mathlib versions;
\item the \lean{PauliCubes} library (about 4{,}700 lines; first written for Lean~4.33.1 and ported with the
edits recorded in the repository), which proves Theorem~\ref{thm:cubes}, Hrube\v{s}'s
Lemma~22 in the form of Lemma~\ref{lem:hrubes}, and the reductions for $S_{4,N}$ and $\mathrm{perm}_{4,N}$;
\item the \lean{IDMM} library (51 files, about 8{,}000 lines) for this paper.
\end{itemize}
The top-level statements are:
{\sloppy
\begin{itemize}
\item \lean{IDMM.id\_bound\_every\_char} (Theorem~\ref{thm:main}): for every field \lean{F : Type u} with
$\mathtt{ringChar}\,F\ne2$, \lean{IDBoundF}, \lean{S4BoundF} and \lean{Perm4BoundF} hold over \lean{F} at
$4173/2500=1.6692$;
\item \lean{IDMM.main\_target : MainTarget}, i.e.\ \lean{IDBound}, \lean{S4Bound} and \lean{Perm4Bound} at
$211321/125000$ over $\C$;
\item \lean{IDMM.main\_target\_almost\_every\_char}: there is a finite \lean{S : Finset Nat} such that for
every field \lean{F : Type u} with $\mathtt{ringChar}\,F\ne2$ and $\mathtt{ringChar}\,F=0\lor\mathtt{ringChar}\,F\notin S$, the same three bounds hold
over \lean{F};
\item \lean{IDMM.main\_target\_char\_zero}: the same bounds over every field of characteristic zero;
\item \lean{PauliCubes.id\_circuit}, \lean{PauliCubes.s4\_circuit} and \lean{PauliCubes.perm4\_circuit}
(Theorem~\ref{thm:cubes}): over every field with $2\ne0$ and for every $\varepsilon>0$, circuits of size
$O(N^{1.727899+\varepsilon})$; \lean{IDMM.cubes\_bound} restates them as \lean{IDBoundF}, \lean{S4BoundF} and
\lean{Perm4BoundF};
\item \lean{IDMM.construction\_almost\_every\_char} and \lean{IDMM.construction\_char\_zero}:
Theorem~\ref{thm:fieldtransfer}(a) and~(b);
\item \lean{IDMM.thin\_main}: the bound $16939/10000$ over $\C$ from the thin endpoint alone;
\item \lean{IDMM.construction\_target : ConditionalTarget} (Theorem~\ref{thm:construction} over $\C$, for every
$\mu\in(\frac12,1)$ and every admissible $w$) and \lean{IDMM.construction\_gen} (any field with $2\ne0$ and
$\sqrt{-1}$, given block oracles);
\item \lean{IDMM.strict\_id\_bound\_every\_char}, \lean{IDMM.strict\_main\_target},
\lean{IDMM.strict\_main\_target\_char\_zero}, \lean{IDMM.strict\_cubes\_bound} and
\lean{IDMM.strict\_main\_target\_almost\_every\_char}: Theorems~\ref{thm:main}, \ref{thm:oairoute} and~\ref{thm:cubes} in the strict
circuit model described in Section~\ref{sec:trusted};
\item \lean{IDMM.two\_mul\_le\_of\_ncSizeLE\_ID}, \lean{IDMM.not\_idBoundF\_of\_lt\_one},
\lean{IDMM.id\_bound\_every\_char\_rat} and \lean{IDMM.id\_bound\_every\_char\_zmod3}: the non-vacuity checks
described there.
\end{itemize}
\par}
Every one of them depends only on Lean's axioms \lean{propext}, \lean{Classical.choice} and \lean{Quot.sound}; no
\lean{sorry}, \lean{native\_decide} or added axiom occurs anywhere.

The repository \url{https://github.com/research-notes-100/identity-polynomial-circuits} contains exactly the
sources verified. On October~29, 2026, from 10:16 to 10:48~p.m.\ Eastern, \texttt{scripts/validate.py} ran on them,
in order: the full \lean{lake build}; the audit modules \lean{IDMM.FinalAudit} and
\lean{IDMM.FinalAuditEveryChar}, which print the signature and axiom list of every theorem above
(\texttt{scripts/audit.py} fails if one is missing); and
\lean{leanchecker --fresh} on both modules, which replays each module and everything it imports through the
kernel in a fresh environment. Every stage exited with status~0. The record in \texttt{validation/release} gives
the SHA-256 hash of every Lean source, pin file and log, and \texttt{scripts/audit.py} checks the repository
against it and checks OpenAI's and Zhou and Wang's files against their upstream hashes. The \texttt{theorem-b}
directory, a separate Lake project on Lean~4.33.1 and Anthropic's formalization, was checked the same day from
4:13 to 4:52~p.m.: a full build, a check that each of its 312 constants uses only the three axioms above, and
kernel replays of all seven of its modules, in place and with \lean{--fresh} (\texttt{validation/theorem-b}).

Zhou and Wang's library and OpenAI's library declare some of the same names (integer and complex versions of shared
modules), so they cannot be imported into one file. The proof of Theorem~\ref{thm:main} therefore uses Zhou and
Wang's library together with the OpenAI modules it needs, but not OpenAI's top-level matrix-multiplication
theorems; the proofs of Theorem~\ref{thm:oairoute} use OpenAI's library and not Zhou and Wang's. Each closure is
built and checked on its own. The import closure of \lean{IDMM.id\_bound\_every\_char} contains 52 modules of Zhou
and Wang's library, 55 of OpenAI's, 5 of harfe's, 22 of \lean{PauliCubes} and 34 of \lean{IDMM}, besides Mathlib.

\subsection{The trusted definitions}\label{sec:trusted}
A formal theorem is only as good as its statement. The statements rest on the following definitions, which we
compared with Hrube\v{s}'s~\cite{Hrubes2026}.
\begin{itemize}
\item \emph{Circuits.} \lean{NCSizeLE f s} holds if some list of at most $s$ elements of the free algebra
$\F\langle V\rangle$, each a variable, a constant of $\F$, or the sum or product (in either order) of two earlier
elements, contains $f$. A list of node values in topological order is the same as a circuit in the sense of
Definition~\ref{def:circuit}, counting inputs as nodes; the only difference is that a gate may use the same earlier
node twice. Read strictly, Definition~\ref{def:circuit} requires the two inputs of a gate to be different nodes.
\lean{StrictNCSizeLE} formalizes that reading by node positions, and \lean{IDMM.ncSizeLE\_strict} turns a circuit
with $s$ nodes into a strict one with at most $2s+1$ nodes: one constant node $1$, and after every node $g$ a copy
$1\cdot g$. Conversely every strict circuit is a circuit (\lean{IDMM.StrictNCSizeLE.ncSizeLE}).
Theorems~\ref{thm:main}, \ref{thm:oairoute} and~\ref{thm:cubes} are proved in the strict model as well.
\item \emph{Non-vacuity.} Every circuit computing $\ID_N$ has at least $2N$ nodes, one for each variable
(\lean{IDMM.two\_mul\_le\_of\_ncSizeLE\_ID}). Hence \lean{IDBoundF F c} is false for every $c<1$, in both models
(\lean{IDMM.not\_idBoundF\_of\_lt\_one}, \lean{IDMM.not\_strictIDBoundF\_of\_lt\_one}): no bound statement holds
with an exponent below $1$. These lower bounds do not show that our exponents are optimal. Theorem~\ref{thm:main} is also instantiated over $\mathbb{Q}$ and over $\mathbb{F}_3$, so
its hypothesis is satisfiable.
\item \emph{Polynomials.} \lean{ID N}, \lean{S4 N} and \lean{perm4 N} are exactly the sums displayed in Section~2.
\item \emph{Fields.} The theorems over general fields quantify over \lean{F : Type u} for an arbitrary universe
\lean{u}.
\item \emph{Bounds.} \lean{IDBound c} (resp.\ \lean{IDBoundF F c}) states: for every $\varepsilon>0$ there is
$C>0$ with \lean{NCSizeLE (ID N)} $\lfloor CN^{c+\varepsilon}\rfloor$ for all $N\ge1$.
\item \emph{Matrix multiplication.} The statement of Theorem~\ref{thm:main} mentions no matrix-multiplication
notion; the input enters only inside the proof, as an integer decomposition checked against Zhou and Wang's
\lean{matrixTensor}, which \lean{IDMM.rankAtMost\_of\_scheme} maps into any commutative ring. The notions of
Section~2.4, used for Theorem~\ref{thm:oairoute}, are OpenAI's definitions, imported unchanged.
\end{itemize}
The statement file of the conditional milestone (\lean{IDMM/Statements.lean}), the circuit definitions and
OpenAI's files are unchanged from the versions in which they were first written.

\section{Discussion}
\begin{enumerate}[label=(\arabic*)]
\item \emph{Every odd characteristic.} Theorem~\ref{thm:main} is uniform across characteristics, with no exceptional
odd primes: one integer decomposition serves every field. (It is an existence proof; neither we nor Zhou and Wang
extract an executable algorithm.) The uncomputed set $S$ of Theorem~\ref{thm:oairoute}(b) remains only in that
route. Lynn~\cite{Lynn2015} studied the analogous question for sum-of-squares formulas over
algebraically closed fields.
\item \emph{Better matrix multiplication feeds in directly.} The exponent is $C(\mu,w)$ with $w$ any rectangular
exponent at aspect $\frac{1-\mu}{\mu}$, and the formal Lemma~\ref{lem:scheme} accepts any rank decomposition. An
improvement of $\omega(1,k,1)$ near $k\approx0.61$, given by one explicit decomposition, would improve
Theorem~\ref{thm:main} with no new circuit argument.
\item \emph{Characteristic two} is excluded by Lemma~\ref{lem:hrubes}, which divides by $2$.
\item \emph{The floor of the cost formula} is $2\Hb(\frac14)=1.6226$, approached only if $\omega=2$. It bounds this
accounting, not circuits: the count is of all words in corner classes, so an identity with fewer nonzero weights, or
evaluation shared across blocks, could go below it, as could a different route from sums of squares to circuits.
\item \emph{Lower bounds.} Nothing here bears on lower bounds; the program of~\cite{HWY2011} needs only
$\sigma(n)\ge n^{1+\varepsilon}$.
\end{enumerate}

\section*{Acknowledgments}
This work rests on the work of many people, and it is a pleasure to thank them.

Pavel Hrube\v{s}, Avi Wigderson and Amir Yehudayoff created the bridge between sums of squares and noncommutative
circuits on which everything here stands, and Hrube\v{s}'s 2024 paper supplied both the first subquadratic bounds
and the exact form of the reduction (his Lemma~22) that we use. Behind them stand Hurwitz and Radon, and Daniel
Shapiro's monograph, which keeps the classical theory accessible. Noam Nisan introduced the noncommutative circuit
model. The algebraic tools are Pascual Jordan and Eugene Wigner's 1928 transformation, Mykhailo Krawtchouk's
polynomials, and the Baur--Strassen theorem.

The main theorem rests on Przemys\l{}aw Uzna\'nski's rectangular bound~\cite{Uznanski2026}, which its author
reports was produced with OpenAI's Codex, and on Yang Zhou and Mingsheng Wang's extension of it to every ring, with
its Lean formalization, released openly~\cite{ZhouWang2026}. Their integer decompositions are what remove every
restriction on the characteristic beyond Hrube\v{s}'s. Their development builds on OpenAI's library, on Jeroen
Zuiddam's work on asymptotic spectra~\cite{Zuiddam2020} and his Lean development
\texttt{omega-nine-quarters-all-fields}~\cite{Zuiddam2026}, and on harfe's formalization of Brouwer's fixed-point theorem~\cite{harfe2026}; the
spectral method goes back to Strassen~\cite{Strassen1988}.

The rectangular matrix-multiplication bounds of the second route, and the $9/4$ bound that started this line, are
due to OpenAI~\cite{OAI2026b,OAI2026a}, who
also formalized them in Lean and released the formalization openly~\cite{OAI2026lean}; we import 509 of their files
unchanged, and without that library the second route would rest on an unverified input. Their work builds on a long line, including Volker Strassen~\cite{Strassen1969,Strassen1973}, Arnold
Sch\"onhage~\cite{Schonhage1981}, Don Coppersmith and Shmuel Winograd~\cite{CoppersmithWinograd1990}, Grazia
Lotti and Francesco Romani~\cite{LottiRomani1983}, Fran\c{c}ois Le~Gall and Florent
Urrutia~\cite{LeGallUrrutia2018}, and Josh Alman, Ran Duan, Virginia Vassilevska Williams, Yinzhan Xu, Zixuan Xu
and Renfei Zhou~\cite{ADVWXXZ2024}; Amir Shpilka's exposition~\cite{Shpilka2026} helped us read it. This project grew out of
studying the October 2026 work of OpenAI and of Anthropic on the Exponential Time Hypothesis. Anthropic's
paper~\cite{Anthropic3SAT} and its Lean-verified \texttt{formal-math} development~\cite{AnthropicFormalMath} led
us to the improved composition bound mentioned in the introduction and from there to this problem; its single-bit
analysis of the crossing tensor anticipates our first Pauli decomposition; and it set the standard of
machine-checked proof that we have tried to meet.

The formal proofs depend on Lean~4, created by Leonardo de~Moura and Sebastian Ullrich and developed by the Lean
FRO and its contributors, and on Mathlib, whose community built, among much else, the model theory of algebraically
closed fields and the Lefschetz principle used in Section~\ref{sec:charp}. We thank A.~Kapila and Vaibhav Tulsyan, whose independent same-day notes~\cite{Kapila2026,Tulsyan2026} on
sums of squares we became aware of during this work.

\paragraph{Use of AI systems.} The mathematics, the Lean proofs and this paper were produced by AI systems under
the supervision of Carson Weitnauer: OpenAI's Codex (in ChatGPT) for research cycles, including the construction of
Sections~\ref{sec:corner}--\ref{sec:assembly}, and for a review of the formalization that contributed
Theorem~\ref{thm:fieldtransfer} and its Lean proof; and Anthropic's Claude for review, the Lean formalization and
the writing of this paper. All theorems and lemmas are machine-checked as described in Section~\ref{sec:lean}.
Carson Weitnauer is responsible for the release of this work. How to attribute authorship of such work is an open
question.

\begin{thebibliography}{99}
\bibitem{ADVWXXZ2024} J.~Alman, R.~Duan, V.~Vassilevska Williams, Y.~Xu, Z.~Xu and R.~Zhou, \emph{More asymmetry
yields faster matrix multiplication}, arXiv:2404.16349, 2024.
\bibitem{ACDM2019} V.~Arvind, A.~Chatterjee, R.~Datta and P.~Mukhopadhyay, \emph{On explicit branching programs
for the rectangular determinant and permanent polynomials}, ISAAC 2019; Chicago J. Theoret. Comput. Sci. 2020.
\bibitem{Anthropic3SAT} Anthropic, \emph{A $2^{o(n)}$ time algorithm for 3-SAT}, October~26, 2026,
\texttt{github.com/anthropics/formal-math}, \texttt{eth/paper.pdf}.
\bibitem{AnthropicFormalMath} Anthropic, \texttt{formal-math} repository, \texttt{eth} development (Lean~4),
commit \texttt{16a5ce4}, 2026.
\bibitem{BaurStrassen1983} W.~Baur and V.~Strassen, \emph{The complexity of partial derivatives}, Theoret.
Comput. Sci. 22 (1983), 317--330.
\bibitem{BCS1997} P.~B\"urgisser, M.~Clausen and M.~A.~Shokrollahi, \emph{Algebraic Complexity Theory}, Springer,
1997.
\bibitem{CoppersmithWinograd1990} D.~Coppersmith and S.~Winograd, \emph{Matrix multiplication via arithmetic
progressions}, J. Symbolic Comput. 9 (1990), 251--280.
\bibitem{harfe2026} harfe, \texttt{fixed-point-theorems-lean4} (Lean~4 formalization of Brouwer's fixed-point
theorem), GitHub, 2026.
\bibitem{Hrubes2024} P.~Hrube\v{s}, \emph{A subquadratic upper bound on sum-of-squares composition formulas},
ECCC TR24-026, 2024; CCC 2024, LIPIcs 300, 12:1--12:11.
\bibitem{Hrubes2026} P.~Hrube\v{s}, \emph{A subquadratic upper bound on Hurwitz's problem and related
non-commutative polynomials}, ECCC TR24-099, 2024; SIAM J. Comput. 55(2) (2026), 451--467.
\bibitem{HWY2011} P.~Hrube\v{s}, A.~Wigderson and A.~Yehudayoff, \emph{Non-commutative circuits and the
sum-of-squares problem}, J. Amer. Math. Soc. 24 (2011), 871--898.
\bibitem{Hurwitz1898} A.~Hurwitz, \emph{\"Uber die Komposition der quadratischen Formen von beliebig vielen
Variabeln}, Nachr. Ges. Wiss. G\"ottingen (1898), 309--316.
\bibitem{JordanWigner1928} P.~Jordan and E.~Wigner, \emph{\"Uber das Paulische \"Aquivalenzverbot}, Z. Physik 47
(1928), 631--651.
\bibitem{Kapila2026} A.~Kapila, \emph{An $N^{\sqrt2+o(1)}$ upper bound for bilinear sum-of-squares composition
formulas}, \texttt{github.com/a-kapila/maths-proofs}, October~27, 2026.
\bibitem{Krawtchouk1929} M.~Krawtchouk, \emph{Sur une g\'en\'eralisation des polyn\^omes d'Hermite}, C. R. Acad.
Sci. Paris 189 (1929), 620--622.
\bibitem{Lean4} L.~de~Moura and S.~Ullrich, \emph{The Lean~4 theorem prover and programming language}, CADE-28,
2021.
\bibitem{LeGallUrrutia2018} F.~Le~Gall and F.~Urrutia, \emph{Improved rectangular matrix multiplication using
powers of the Coppersmith--Winograd tensor}, SODA 2018.
\bibitem{LottiRomani1983} G.~Lotti and F.~Romani, \emph{On the asymptotic complexity of rectangular matrix
multiplication}, Theoret. Comput. Sci. 23 (1983), 171--185.
\bibitem{Lynn2015} B.~Lynn, \emph{Sums-of-squares formulas over algebraically closed fields},
arXiv:1510.05369, 2015.
\bibitem{MacWilliamsSloane1977} F.~J.~MacWilliams and N.~J.~A.~Sloane, \emph{The Theory of Error-Correcting
Codes}, North-Holland, 1977.
\bibitem{Mathlib2020} The mathlib Community, \emph{The Lean mathematical library}, CPP 2020.
\bibitem{Nisan1991} N.~Nisan, \emph{Lower bounds for non-commutative computation}, STOC 1991, 410--418.
\bibitem{OAI2026a} OpenAI, \emph{An upper bound of $9/4$ for the matrix multiplication exponent}, preprint,
October~2, 2026, \texttt{github.com/openai/math}.
\bibitem{OAI2026b} OpenAI, \emph{Complex matrix multiplication below $2.258$ and rectangular bounds}, preprint,
September~24, 2026, \texttt{github.com/openai/math}.
\bibitem{OAI2026lean} OpenAI, Lean formalization of the matrix-multiplication bounds,
\texttt{github.com/openai/math}, \texttt{lean/OAI/LinearAlgebra/MatrixMultiplication}, commit \texttt{7ef5c0c},
2026.
\bibitem{Radon1922} J.~Radon, \emph{Lineare Scharen orthogonaler Matrizen}, Abh. Math. Sem. Univ. Hamburg 1
(1922), 1--14.
\bibitem{Schonhage1981} A.~Sch\"onhage, \emph{Partial and total matrix multiplication}, SIAM J. Comput. 10 (1981),
434--455.
\bibitem{Shapiro2000} D.~B.~Shapiro, \emph{Compositions of Quadratic Forms}, de Gruyter, 2000.
\bibitem{Shpilka2026} A.~Shpilka, \emph{Exposition of OpenAI's matrix multiplication in exponent $9/4$ (and an
explicit algorithm)}, ECCC TR26-281, 2026.
\bibitem{Strassen1969} V.~Strassen, \emph{Gaussian elimination is not optimal}, Numer. Math. 13 (1969), 354--356.
\bibitem{Strassen1988} V.~Strassen, \emph{The asymptotic spectrum of tensors}, J. Reine Angew. Math. 384
(1988), 102--152.
\bibitem{Strassen1973} V.~Strassen, \emph{Vermeidung von Divisionen}, J. Reine Angew. Math. 264 (1973), 184--202.
\bibitem{Uznanski2026} P.~Uzna\'nski, \emph{Rectangular matrix multiplication from polynomial products},
arXiv:2610.10494v2, October~11, 2026 (v1, October~7, 2026, titled \emph{Rectangular matrix multiplication from
shared-leg entropy}).
\bibitem{ZhouWang2026} Y.~Zhou and M.~Wang, \emph{Shared-leg rectangular matrix multiplication over arbitrary
rings}, research note and Lean~4 formalization, \texttt{github.com/Hugin-Hilbert/rectangular-matmul-rings},
October~14, 2026, commit \texttt{285f513}.
\bibitem{Zuiddam2020} J.~Zuiddam, \emph{Algebraic complexity, asymptotic spectra and entanglement polytopes},
PhD dissertation, University of Amsterdam, 2018; updated version, 2020.
\bibitem{Zuiddam2026} J.~Zuiddam, \texttt{omega-nine-quarters-all-fields} (Lean~4 development), GitHub, 2026.
\bibitem{Tulsyan2026} V.~Tulsyan, \texttt{eth-hamiltonicity} repository, \texttt{research/pieces.md}, October 27,
2026.
\end{thebibliography}

\end{document}
